\section{Related Work}

\noindent\textbf{Contextual and co-occurrence bias in recognition models.} Reliance on co-occurring context as a predictive shortcut is well-documented for older architectures, but less so for the modern stack. \emph{Causal interventions} demonstrate the dependence by manipulating features or inputs: via causal inpainting on convolutional neural network (CNN) features~\cite{shetty2019not}, and via feature-space splitting on COCO-Stuff~\cite{caesar2018coco,singh2020don}. 
\emph{Out-of-context benchmarks} measure model behavior when objects appear in atypical scenes~\cite{acharya2022detecting,bomatter2021pigs}. \emph{Concept-graph approaches} discover and mitigate co-occurrence biases via concept-graph clique sampling~\cite{chakraborty2024visual} or open-set foundation-model tagging~\cite{sarridis2025mavias}.
More recent work has examined context bias in modern models, but only in narrower settings: XAI for segmentation and detection~\cite{dreyer2023revealing}, nuisance-based robustness benchmarks~\cite{zhao2022ood}, referring expression generation~\cite{merlo2026cooco}, and compositional action recognition~\cite{ahn2026can}.
In contrast to these narrower studies, we apply one diagnostic across five model families, complemented by \textbf{STRANGE-Bench}, which covers co-occurrence regimes absent from natural data, and we also test mitigations.

\begin{table}[t!]
\centering
\caption{\textbf{Taxonomy of popular data augmentations for visual recognition.} None are designed to break co-occurrence shortcuts. Columns: Std (standard: resize, crop, flip), MixUp~\cite{zhang2017mixup}, Mosaic~\cite{bochkovskiy2020yolov4}, YOLOX~\cite{ge2021yolox} (MixUp + Mosaic), SCP (Simple Copy-Paste~\cite{ghiasi2021simple}), XPaste~\cite{zhao2023x}, InstAug (InstanceAug~\cite{kupyn2024dataset}), MFusion (MosaicFusion~\cite{xie2025mosaicfusion}). Static augmentations operate within a single image; dynamic ones introduce new instances at training time.}
\label{tab:augmentations}
\small
\begin{tabular}{@{}lcccccccc@{}}
\toprule
\textbf{Property} & \textbf{Std} & \textbf{MixUp} & \textbf{Mosaic} & \textbf{YOLOX} & \textbf{SCP} & \textbf{XPaste} & \textbf{InstAug} & \textbf{MFusion} \\
\midrule
Non-generative         & \cmark       & \cmark         & \cmark         & \cmark          & \cmark       & \xmark          & \xmark           & \xmark           \\
Generative             & \xmark       & \xmark         & \xmark         & \xmark          & \xmark       & \cmark          & \cmark           & \cmark           \\
\midrule
Static                 & \cmark       & \xmark         & \xmark         & \xmark          & \xmark       & \xmark          & \cmark           & \xmark           \\
Dynamic                & \xmark       & \cmark         & \cmark         & \cmark          & \cmark       & \cmark          & \xmark           & \cmark           \\
\midrule
\textbf{Context-aware} & \xmark       & \xmark         & \xmark         & \xmark          & \xmark       & \xmark          & \xmark           & \xmark           \\
\bottomrule
\end{tabular}
\vspace{-10pt}
\end{table}


\noindent\textbf{Data augmentation for visual recognition.} Data augmentation shapes the training distribution that models see, but existing methods were not designed with co-occurrence shortcuts in mind. \emph{Non-generative methods} perturb individual images (flip, crop, color jitter) or stitch existing instances together (MixUp~\cite{zhang2017mixup}, Mosaic~\cite{bochkovskiy2020yolov4}, Simple Copy-Paste~\cite{ghiasi2021simple}). \emph{Generative methods} synthesize new instances using text-to-image diffusion (XPaste~\cite{zhao2023x}, InstanceAug~\cite{kupyn2024dataset}, MosaicFusion~\cite{xie2025mosaicfusion}). Across both families (see Table~\ref{tab:augmentations}), pasted or generated instances are sampled from the natural co-occurrence distribution. Whether this preserves typicality bias, and whether targeted modifications can reduce it, have not been systematically studied. We present such a study on representative model families, and
propose \textbf{Anchor-Atypicality Sampling (AAS)}, a sampling rule that plugs into any paste-based augmentation and reduces typicality bias while preserving overall performance.

\noindent\textbf{Modern visual recognition models.} The vision recognition stack has evolved rapidly along architectural and supervision axes. \emph{Specialist detection transformers} such as DETR~\cite{carion2020end}, Deformable DETR~\cite{zhu2021deformable}, and Co-DETR/Co-DINO~\cite{zong2023detrs} formulate detection end-to-end via (deformable) attention. \emph{Open-vocabulary detectors} such as GLIP~\cite{li2022grounded} and Grounding DINO~\cite{liu2024grounding} extend detection to text-conditioned localization. \emph{Self-supervised foundation models} such as DINOv2~\cite{oquab2024dinov} and DINOv3~\cite{simeoni2025dinov3} learn general-purpose visual representations from hundreds of millions of unlabeled images. \emph{Vision--language models} (CLIP~\cite{radford2021learning}, OpenCLIP~\cite{cherti2023reproducible}) and \emph{large multimodal models} (LLaVA-OneVision~\cite{li2024llava}) use web-scale image--text pretraining for zero-shot recognition and image understanding. Each of these design choices, from deformable attention focusing on object regions to language grounding and web-scale data diversity, is expected to be resistant to context shortcuts. Yet none is trained or evaluated with such shortcuts in mind, leaving their susceptibility to typicality bias underexamined, a gap our diagnostic, benchmark, and mitigation directly address.
